\documentclass[11pt]{article}

\usepackage[margin=2.4cm]{geometry}
\usepackage[T1]{fontenc}
\usepackage{lmodern}
\usepackage{amsmath,amssymb}
\usepackage{xcolor}
\usepackage{enumitem}
\usepackage{titlesec}
\usepackage{hyperref}

\definecolor{hgpurple}{RGB}{106,61,154}

\titleformat{\section}{\Large\bfseries\color{hgpurple}}{\thesection}{0.6em}{}
\titleformat{\subsection}{\large\bfseries\color{hgpurple}}{}{0em}{}
\hypersetup{colorlinks=true, linkcolor=hgpurple, urlcolor=hgpurple}

\title{\color{hgpurple}\textbf{Hamiltonian-Geometric Optimization}\\[2pt]
  \large Speaker Script and Anticipated Q\&A (14-Slide Deck)}
\author{Sparsho Chakraborty et al. --- Capacitor Fringes Project}
\date{\today}

\begin{document}
\maketitle

\noindent This document accompanies \texttt{hamiltonian\_geometric\_short\_presentation.tex}.
Part~\ref{sec:script} is a read-aloud script, one block per slide (each roughly 30--45 seconds,
so the full talk runs about 8--10 minutes). Part~\ref{sec:qa} is a set of hard, likely audience
questions with grounded answers, meant to be read \emph{after} the talk, in preparation for
questions --- not to be presented.

\section{Speaker Script}
\label{sec:script}

\subsection*{Slide 1 --- Title}
``Today I'm presenting the Hamiltonian-geometric optimization framework --- a way of designing
optimizers by treating training itself as a physical system. This is a short walkthrough: the
math, how we validated it, and what the benchmarks actually show.''

\subsection*{Slide 2 --- Core Question}
``The question we're testing is simple to state: can one Hamiltonian-geometric update rule work
across very different settings --- classical ML, LLM-style training, physics-informed networks,
and quantum control? We treat the model's parameters as coordinates, the loss as potential
energy, and the optimizer's momentum as a physical momentum moving on a curved metric manifold.
And I want to be upfront about the claim: it's bounded. We are not claiming this beats Adam
everywhere --- only that it's strongest when the problem actually has exploitable curvature.''

\subsection*{Slide 3 --- Methodology Flowchart}
``This is the whole pipeline in one picture, so you don't have to hold every equation in your
head at once. We start from the parameters and the loss, build a metric from the regularized
Hessian, plug that into a Hamiltonian, and get the equations of motion out of it. Three
correction forces --- geometric, memory, and spectral-entropy --- feed into one generalized
update rule, which also collapses back down to familiar optimizers like SGD or Adam as special
cases. Everything below this box is validated first, analytically, before we ever look at a
benchmark curve.''

\subsection*{Slide 4 --- Hamiltonian Formulation}
``Here's the actual Hamiltonian: kinetic term is momentum contracted through the inverse metric,
potential term is the loss. Differentiating gives the equations of motion, and the interesting
part is the second term in the momentum equation --- that's the geometric force, and it only
exists because the metric depends on where you are in parameter space. We add dissipation on
top, because plain Hamiltonian mechanics conserves energy, and we want the loss to actually go
down, not oscillate forever.''

\subsection*{Slide 5 --- One Update Rule, Many Optimizers}
``The point of this slide is that we didn't invent a totally separate optimizer family. SGD and
heavy-ball fall out when the metric is flat. Adam and RMSProp fall out when you keep only the
diagonal of the metric. Natural gradient is the full metric without the momentum and memory
terms. Our Hamiltonian-geometric optimizer is what you get when you keep everything. And modern
optimizers like Shampoo, K-FAC, or Muon all sit somewhere on that same spectrum, from a full
matrix down to a single scalar.''

\subsection*{Slide 6 --- Analytical Validation}
``This slide answers the question people always ask: how do you know these equations are
actually correct? Every equation here traces back to a specific reference: my own
literature-review PDF, which is the design document for this optimizer's architecture. We didn't
just implement it and hope --- we wrote a separate numerical validation module,
\texttt{paper\_math\_validation.py}, that checks each equation with finite differences on a
fixed, seeded toy problem where the metric genuinely curves with the parameters. Three kinds of
checks: `literal' means we tested an equation exactly as the PDF prints it. `Completed' means
the PDF states something but never actually derives it --- like the gradient of the spectral
entropy --- so we derived a specific closed form ourselves and checked it against a numerical
gradient. `Counterexample' is the honest one: the PDF's condition for the Hessian regularization
doesn't actually guarantee a positive-definite metric, we show that failure on purpose, and then
confirm our shipped code fix repairs it, over hundreds of random trials. All seven checks pass to
numerical precision.''

\subsection*{Slide 7 --- Benchmark Evidence Snapshot}
``Four results at once here, on purpose, so you can see the workload variety in one glance:
MNIST accuracy, an AlgoPerf-style validation loss curve, DeepOBS-style validation accuracy, and
an industry-style LLM run on WikiText-2. Across these, the Hamiltonian-geometric optimizer posts
the best or near-best validation loss and accuracy. But these are small-scale runs --- I'd call
this evidence of correct mechanism, not a scaling proof. The next slide puts exact numbers on all
of this.''

\subsection*{Slide 8 --- Benchmark Results Table (Classical, PINN \& LLM)}
``So here are the actual numbers behind those curves. On MNIST, AlgoPerf-style, DeepOBS-style,
and the capacitor PINN, Hamiltonian-geometric reports the best validation or test loss of
everything we compared against. The LLM row is the one I want to be honest about: on WikiText-2,
AdaFactor actually edges us out slightly on validation loss --- 6.765 versus our 6.778. But we
still clearly beat AdamW and Lion, and we finish fastest of the five in wall-clock time.
Muon-lite diverges entirely here, which is a reminder that these are untuned, single-seed runs,
not a scaling claim.''

\subsection*{Slide 9 --- Modern Optimizer Comparison}
``Now the harder comparison: against modern optimizers people actually use today. Left panel is
compute cost versus accuracy --- Adam, our method, and Muon lead the pack. Middle panel is a
10-seed sweep, and the near-tie holds up under that noise. Right panel is where it gets more
honest: as we scale model width, Adam scales a bit more consistently than we do. So: competitive,
genuinely, but not a decisive win. And to be transparent, the version benchmarked here is a
diagonal-metric reduction, not the full curvature-memory-entropy architecture.''

\subsection*{Slide 10 --- Benchmark Results Table (Modern Optimizers)}
``And the numbers behind that comparison. Ranked by best-epoch validation loss on the spiral
MLP: Adam, Hamiltonian-geometric, and Muon are within four ten-thousandths of each other ---
that's noise, not a real ordering. AdamW and Lion trail by a clearer margin, and Shampoo is both
slower per epoch and slightly worse here. The honest read is a narrow leading group of three, not
a win for any one of them.''

\subsection*{Slide 11 --- Extended Ablation Study: Design \& Evidence}
``Everything up to here used the same seeds and setups as the original report. This next part is
a separate, more rigorous replication we ran afterward, specifically to stress-test those claims:
a full $2^3$ factorial crossing the geometric, memory, and spectral forces, with paired seeds ---
40 seeds over 300 steps on the architecture task, 10 seeds on the chaotic quantum task, and 10
seeds on a CUDA scaling task using the diagonal reduction. The point of showing three panels at
once is that this is one coherent study, not three unrelated plots.''

\subsection*{Slide 12 --- Extended Ablation Study: Key Numbers}
``And here's what that rigor actually found. On the architecture task, it confirms the earlier
story precisely: the spectral term only has a significant effect --- a large one, $p$ under
$10^{-4}$ --- when the metric is the genuine Hessian, not an unrelated control metric. But on the
chaotic quantum task, I want to be direct: the result is negative. Turning on the geometric force
significantly \emph{lowers} fidelity, and every simple baseline we tested --- heavy-ball, entropy
descent, AdamW --- beats every full Hamiltonian-geometric variant, while costing an order of
magnitude less runtime. Same story on the CUDA scaling ablation: the diagonal reduction trails
plain SGD and AdamW there too. I'd rather show you this than hide it --- it's exactly the kind of
rigor that should update the claims, not just decorate them.''

\subsection*{Slide 13 --- Limitations}
``Four honest caveats. Scale: everything here is small; a full n-by-n metric is a mechanism
demonstration, not something we're claiming replaces Adam in production. Statistics: the earlier
classical and LLM benchmarks are mostly single-seed, but the modern-optimizer sweep and the
extended ablation now give us real multi-seed variance --- and that's precisely where the
negative quantum and GPU results show up. Tuning: learning rates come from modest grids, not
exhaustive search. And comparability: both the modern-optimizer run and the CUDA ablation use a
diagonal reduction, not the full curvature-memory-entropy architecture.''

\subsection*{Slide 14 --- Summary and Thank You}
``Four takeaways. The formulation holds together once you're explicit about the spectral-force
sign and the memory-metric construction. The empirical evidence is strongest exactly where the
theory predicts it should be --- where there's real, exploitable curvature. Against modern
optimizers, we're in the leading group, not decisively ahead of Adam or Muon. And the extended
ablation sharpens that picture rather than softening it: it's a real, honest positive on the
spectral term under a true Hessian metric, and a real, honest negative on quantum control and
GPU-scale diagonal reduction. Happy to take questions, especially on the validation methodology,
the ablation results, or where we'd want to push this next.''

\clearpage
\section{Anticipated Questions and Answers}
\label{sec:qa}

\noindent The questions below are ordered roughly hardest-and-most-important first. Each answer
is grounded in a specific number or result already in the codebase (the numeric validation module
and the consolidated report), so they can be defended on the spot rather than improvised.

\begin{enumerate}[leftmargin=1.4em, itemsep=0.9em]

\item \textbf{How do you actually know these equations are correct, rather than just plausible-looking physics?}\\
We don't validate against the benchmark tasks themselves --- we validate the equations directly,
symbol by symbol, in \texttt{src/paper\_math\_validation.py}. Each of the seven checks compares a
closed-form expression from the reference PDF against an independent finite-difference
computation on a fixed, seeded synthetic problem, to machine precision (errors of $10^{-9}$ to
$10^{-16}$). Critically, this toy problem was built specifically so the metric $g(\theta)$
depends on $\theta$; the project's own capacitor benchmark has a metric that is \emph{constant}
in $\theta$, which would make $F_{\mathrm{geo}}$ and $\nabla_\theta S(g)$ trivially zero and
prove nothing.

\item \textbf{Is the improvement over Adam statistically significant, or just noise?}\\
Borderline, and we say so explicitly. On the 10-seed modern-optimizer sweep, a paired $t$-test
(licensed because every optimizer trains on the same paired $\theta_0$/dataset draw per seed)
gives HG vs.\ Adam: $t(9){=}{-2.234}$, $p{=}0.052$ --- borderline, HG lower on 7 of 10 seeds. HG
vs.\ Muon: $t{=}{-0.838}$, $p{=}0.424$ (not significant). HG vs.\ AdamW: $t{=}{-1.457}$,
$p{=}0.179$ (not significant). $n=10$ is underpowered to settle this either way; the honest claim
is a weak directional signal against Adam specifically, and no signal against Muon or AdamW.

\item \textbf{The memory force is supposed to help, so why does it hurt on the quantum benchmark?}\\
On the unsafeguarded kicked-top ablation, turning on the memory force alone moves final loss from
$0.0491$ (the $\beta{=}0.7$ default) to $0.0631$ --- worse. The explicit geometric force is worse
still ($0.0491\to0.1949$) and over $10\times$ slower, because both differentiate a metric that is
itself built from finite differences. Only the spectral-entropy force helps on this task
($0.0491\to0.0184$). The honest reading: momentum tuning dominates on this benchmark, and the
three correction terms have a \emph{mixed}, task-dependent effect, not a uniformly positive one.
Separately, on the PINN ablation, the picture flips: under a toy (parameter-independent) metric
only the memory term is significant ($p{=}7.4\times10^{-4}$, a small $\sim$12\% effect), but under
the true Hessian metric the spectral term becomes highly significant
($p{=}3.2\times10^{-8}$) while memory and geometric forces are not. The takeaway we lead with is
that these terms matter most --- and differently --- depending on whether the metric is actually
tied to real loss curvature. The extended, paired-seed ablation (10 seeds) sharpens this further
in the geometric force's disfavor specifically: $G$ has a significant \emph{negative} effect on
fidelity ($t$-effect $0.88$, $p{=}1.57\times10^{-4}$), and every full-HG variant (fidelity
$0.86$--$0.98$) is beaten by every simple baseline tested --- heavy-ball ($0.993$), entropy
descent ($0.988$), AdamW ($0.986$) --- at $10$--$20\times$ lower runtime.

\item \textbf{Your own extended ablation study shows Hamiltonian-geometric losing on quantum control and GPU scaling. Doesn't that undercut the whole pitch?}\\
It sharpens the pitch rather than undercutting it, and we'd rather present it than bury it. The
extended ablation was commissioned specifically to stress-test the original claims with more
seeds and a full factorial design, and it does two things at once: it \emph{strengthens} the
architecture claim (spectral term significant only under the true Hessian metric, $p{=}9.25\times
10^{-5}$, replicated at 40 seeds instead of the smaller earlier run), and it \emph{reverses} the
naive expectation on quantum control and GPU-scale diagonal reduction, where every HG variant
loses to simple baselines. The honest overall claim was always "strongest where full-metric
curvature is exploitable, not uniformly superior" --- this result is that claim being confirmed
in both directions, including the direction that doesn't flatter us.

\item \textbf{On WikiText-2, AdaFactor beat you. Doesn't that undercut the "works across domains" claim?}\\
Yes, on that specific metric, and we report it rather than hide it: AdaFactor's validation loss
is $6.765$ against our $6.778$. But it's one 200-step, single-seed, untuned run --- not a scaling
claim. In the same run we clearly beat AdamW ($6.843$) and Lion ($7.032$), and finish fastest in
wall-clock time; Muon-lite diverges outright (loss $9.31$), which itself shows how sensitive these
small, untuned comparisons are. The claim we defend is "competitive across domains," not "best in
every domain."

\item \textbf{Isn't "lambda > 0 guarantees positive-definiteness" obviously true? Why call that a bug?}\\
It's a real counterexample, not pedantry: for the indefinite Hessian
$\mathrm{diag}(-5, 0.1, 3)$ and the PDF's literal $\lambda{=}10^{-3}$, the regularized matrix
$H+\lambda I$ still has a negative eigenvalue, because $\lambda$ is far smaller than the most
negative eigenvalue's magnitude. A fixed constant can never work for every Hessian; the
regularization has to depend on the matrix's own spectrum. The project's shipped fix,
\texttt{src.pinn.\_positive\_definite\_metric}, does exactly that, and we verify it produces a
strictly positive-definite metric on that counterexample and across 200 random indefinite trials.

\item \textbf{You add a dissipation term, but Hamiltonian systems conserve energy. Isn't that a contradiction?}\\
It's intentional, not a contradiction we've glossed over: pure Hamiltonian flow conserves $H$
exactly, which is useless for optimization --- we need the loss to decrease. The Rayleigh
dissipation term is deliberately added on top of the conservative dynamics, making the overall
system a damped (port-Hamiltonian-style) system rather than an energy-conserving one. The
"Hamiltonian" in the name describes where the structure (metric, momentum, curvature) comes from,
not a claim that the implemented dynamics conserve energy.

\item \textbf{The modern-optimizer comparison uses a diagonal reduction of your method --- isn't that just comparing Adam to Adam?}\\
Largely, yes, and the report says this outright: by our own special-case correspondence, that
diagonal reduction is already "a diagonal-metric approximation to the Hamiltonian optimizer,"
i.e.\ close in spirit to Adam with a different momentum discretization. The near-tie there is
expected, not evidence of the framework's distinctive machinery beating Adam. The full machinery
--- $F_{\mathrm{geo}}$, $F_{\mathrm{mem}}$, the spectral term together --- is only exercised on
the PINN (exact Hessian metric) and quantum-control benchmarks.

\item \textbf{Could this scale to a billion-parameter model?}\\
Not as implemented, and we don't claim it could. A full $n\times n$ metric costs $O(n^2)$ memory
and roughly cubic cost to invert or factor, which is infeasible past small networks --- exactly
why Shampoo and K-FAC exist, using Kronecker-factored approximations to cut that cost (e.g.\
$O((mn)^3)\to O(m^3+n^3)$ for a layer shaped $m\times n$). Getting a cheap structured
approximation of our full metric, the way Shampoo/K-FAC do, is future work, not something we've
benchmarked.

\item \textbf{How is this actually different from just using natural gradient, K-FAC, or Shampoo?}\\
Those methods are different points on the \emph{same} metric-structure hierarchy we use to
classify optimizers --- specific approximations to the same underlying Hamiltonian object (full
metric $\to$ Kronecker-factored $\to$ diagonal $\to$ scalar). What none of them include is our
other two ingredients: the non-Markovian memory force built from an exponentially-weighted
gradient history, and the spectral-entropy regularization on the metric's own eigenvalue
distribution. The claim is that curvature alone (natural gradient/K-FAC/Shampoo) is one piece;
momentum-as-phase-space-dynamics and the two correction forces are what's distinctive here.

\item \textbf{What result would actually change your mind about this framework?}\\
Two things would: if, under a genuine full-metric setting (the exact-Hessian PINN or
quantum-control benchmarks) with many more seeds, the additive correction terms showed no
significant, reproducible effect at all --- rather than the mixed, sign-flipping effect we
currently see --- that would undercut the core mechanistic claim. And if no cheap structured
approximation to the full metric could be found, that would cap the framework's practical
relevance at small-scale, mechanism-level demonstrations rather than anything deployable.

\end{enumerate}

\end{document}
